\section{Monitoring agents against specifications (WF3)}
\label{sec:specs}

Workflow \textbf{(WF3) Monitoring} is where the STL machinery of RLRom first
enters the picture. Using the cart-pole running example, the section walks
through the two ingredients it requires, both prepared declaratively in the
configuration: \emph{interfacing the signals} (mapping the raw observation and
action vector onto named quantities the formulas can refer to), and \emph{writing
very simple specifications} that formalize the \emph{original} goal of the
environment. We then monitor these specifications against an agent that was
\emph{not} trained with them --- the HuggingFace pretrained PPO of
Section~\ref{sec:testing} --- using \texttt{cfg0\_hug\_specs.yml}, which
\emph{passively} monitors the agent and reports robustness-based metrics without
shaping any reward.

Throughout the rest of this chapter we walk through the four workflows of
Section~\ref{sec:workflows} on the running example used in the remaining
sections: the classic \textit{CartPole-v1} Gymnasium environment \cite{gymnasium}.
Its state is a $4$-dimensional vector
$(x,\,\dot x,\,\theta,\,\dot\theta)$ --- cart position, cart velocity, pole angle
and pole angular velocity --- and the action is a push left or right. The episode
terminates once the pole falls (angle beyond a threshold) or the cart leaves the
track. Because each workflow builds on the previous one --- from simple testing,
through training, to STL-conditioned reward machines --- the same
environment and configuration files illustrate how the library supports tasks of
increasing complexity.

To talk about the task in STL, we first \emph{declare signals} that map the raw
observation and action onto named quantities. In \textsc{RLROM} this is done in the
\texttt{cfg\_specs} block through \texttt{action\_names} and \texttt{obs\_names}:
each entry pairs a signal name with a Python expression evaluated at every step
(Listing~\ref{lst:signals}).

\begin{cfglisting}[caption={Signal declarations from the CartPole observation.},label={lst:signals}]
cfg_specs:
  action_names:
    push: "action"
  obs_names:
    x:         "obs[0]"
    x_dot:     "obs[1]"
    theta:     "obs[2]"
    theta_dot: "obs[3]"
\end{cfglisting}

\noindent
The key \texttt{specs} points to a standalone \texttt{.stl} file (here
\texttt{cartpole.stl}, Listing~\ref{lst:stl}) that contains the actual formulas.
Keeping the formulas in a dedicated file, separate from the usage configuration,
lets the same specification be shared by several runs (testing, training, different
roles) --- a central property of the design. The file begins by declaring the
signal set --- which must be consistent with the \texttt{obs\_names}/
\texttt{action\_names} in the configuration --- and then defines a small library of
formulas.

\begin{stllisting}[caption={STL specification for CartPole
(\texttt{examples/cartpole/cartpole.stl}).},label={lst:stl}]
signal push, x, x_dot, theta, theta_dot, reward

param epsi = 0.03
pole_up          := abs(theta[t]) < epsi
keep_up          := alw_[0, 100] pole_up
phi_goal_keep_up := ev_[0, 400] keep_up

(...)

\end{stllisting}

\noindent
A few formulas give the flavor of the language.
\begin{itemize}
  \item \texttt{pole\_up} is a simple instantaneous predicate (the pole is within
  $\varepsilon$ of vertical), from which \texttt{keep\_up} = ``remain balanced for
  the next $100$~steps'' and \texttt{phi\_goal\_keep\_up} = ``eventually, within
  $400$~steps, stay balanced for $100$~steps'' are derived by nesting temporal
  operators.
  \item \texttt{cart\_in\_A} defines a narrow target region of the track
  ($x\in[-0.7,-0.5]$), and \texttt{stay\_A} = ``remain in $A$ for the next
  $50$~steps''; \texttt{not\_A} is its negation. These support the sequential
  reward machine of Section~\ref{sec:rewardmachines}.
  \item \texttt{phi\_left\_goal} / \texttt{phi\_right\_goal} express directional
  objectives --- ``eventually, for $100$~steps, keep the cart in the left (resp.\
  right) half of the track''.
  \item \texttt{stuck\_left} states an inertness property (never leave the right
  region), useful as a \emph{negative} or monitoring-only formula.
\end{itemize}

Several signals (\texttt{push}, \texttt{reward}, \texttt{theta\_dot}) are declared
but not used by these formulas; they are available to any formula and to other
configs loading the same file. Declaring the full signal set once, in the shared
specification, keeps the various role-based usages consistent.

Listing~\ref{lst:hug} shows \texttt{cfg0\_hug\_specs.yml}, which loads a
\emph{pretrained} PPO checkpoint for CartPole from HuggingFace, declares the
signals, points to \texttt{cartpole.stl}, and selects the goal formula
\texttt{phi\_goal\_keep\_up} as an \emph{eval formula} --- a formula monitored
for metrics only, and not used to shape any reward.

\begin{cfglisting}[caption={Testing a HuggingFace pretrained PPO on CartPole
against an STL formula (\texttt{cfg0\_hug\_specs.yml}).},label={lst:hug}]
env_name: CartPole-v1
model_path: huggingface
model_name: sb3/ppo-CartPole-v1

cfg_specs:
  debug_formulas: False
  debug_signals: False
  action_names:
    push: "action"
  obs_names:
    x:     "obs[0]"
    x_dot: "obs[1]"
    theta: "obs[2]"
    theta_dot: "obs[3]"
  specs: cartpole.stl
  eval_formulas:
    phi_goal_keep_up:

cfg_test:
  init_seeds: 1
  num_ep: 10
  num_steps: 100
  render_mode: 'human'
\end{cfglisting}

\noindent
The run is launched with
\begin{center}\ttfamily rlr test examples/cartpole/cfg0\_hug\_specs.yml
\end{center}
which plays the requested episodes while the \texttt{SpecWrapper} monitors the
selected formulas online. During a step the wrapper:
\begin{enumerate}
  \item samples the declared signals from the observation/action;
  \item updates the STL monitor with the new sample;
  \item reads the robustness of each formula and, where relevant, of its lower/upper
  bound;
  \item stores the per-step robustness trace for later analysis.
\end{enumerate}

Across the episodes the tester aggregates, per formula, metrics such as the
per-step robustness trace, the fraction of steps where the formula holds, and ---
for interval/temporal formulas --- whether the goal was eventually reached (initial
satisfaction and initial robustness). These robustness-based metrics provide a much
more interpretable report card than a raw scalar reward: one can read directly
whether ``the pole stayed up within the specified horizon and margin,'' and by what
amount, rather than inferring it from a score.

\paragraph{Remark.} Because the environment is wrapped by \texttt{SpecWrapper}, the
monitoring is \emph{online}: robustness is computed incrementally during the
episode, so the same wrapper can later be plugged into training
(Section~\ref{sec:training}) and reward machines (Section~\ref{sec:rewardmachines})
without any change to the monitoring logic.

%%% Local Variables:
%%% mode: LaTeX
%%% TeX-master: "../rlrom_chapter"
%%% End:
